\section{The Fathers: Life That Does Not Arise from Us}

The Fathers begin elsewhere. They do not extend Aristotle's analysis of motion until it becomes creation. On the ground of Christian revelation, they confess that what is made begins and continues under the will of God.

\subsection{Begun and continued by another}

Irenaeus is answering the claim that a soul must be unbegotten in order to endure. He distinguishes God, who alone is without beginning, from everything made. Created things begin and continue as long as God wills them to exist. The soul is not life itself but participates in life bestowed by God; ``life does not arise from us.''\endnote{Irenaeus, \emph{Against Heresies} II.34.2--4, in Roberts and Rambaut's public-domain \emph{Ante-Nicene Fathers} (ANF) English. His immediate subject is the soul's duration, not a later scholastic theory of \latin{esse}; the wording is used only for the asymmetry it directly states.} The asymmetry is plain: continuance is conferred; the giver does not receive continuance from the recipient.

Athanasius sharpens the Creator--creature distinction against both spontaneous origin and formation from coeternal matter. In \emph{On the Incarnation}, God makes the universe from nothing through the Word and grants existence out of goodness. In \emph{Against the Heathen}, the Word also orders and preserves the created whole. Human beings are mutable creatures; incorruption is not native to humanity but a gift. The governing judgment remains: the Word is not one fragile participant among others. He is the divine giver through whom creatures are called into being and held in order.\endnote{Athanasius, \emph{On the Incarnation} 3--5 and \emph{Against the Heathen} 40--42, in the public-domain \emph{Nicene and Post-Nicene Fathers} (NPNF) working English. His language of corruption and return toward non-being also serves the soteriological account of sin and redemption; it is not flattened here into a single scholastic formula.}

\subsection{More inward than the self}

Augustine turns the same dependence inward. At the opening of the \emph{Confessions} he asks how he could exist unless God were in him, then corrects the direction: rather, unless he were in God. He later calls God more inward than his inmost self and higher than his highest part. This is not pantheism. In Book VII he rejects the thought that God is a spatial substance spread through the world, then apprehends the immutable Light above his mutable mind. Divine transcendence and divine intimacy intensify together: the cause of being is not a part enclosed within the creature, yet no created cause is present at a deeper explanatory level.

For Augustine, dependence is already present in the mind's act of judgment. His changeable mind judges mutable things by unchangeable truth; the life by which he seeks God is upheld in the act of seeking.\endnote{Augustine, \emph{Confessions} I.2.2; III.6.11; VII.1.1--2, 10.16--11.17, and 17.23, in the public-domain NPNF working English. These loci govern the reversal of spatial assumption, the union of inwardness and transcendence, and the mutable mind's judgment by unchangeable truth.}

Across these witnesses, one uncreated God calls creatures from nothing through the Word; created life and duration are gift, and divine presence does not dissolve the distinction between giver and recipient. The First is no longer only the terminus of motion. He is the giver from whom created life begins and by whom it continues.

The resulting asymmetry requires a more exact account of relation: created life depends really upon the giver, while the giver does not receive life or continuance from what is given.
